\documentclass{article}
\usepackage[UTF8]{ctex}
\usepackage{amsmath}
\usepackage{amssymb}
\usepackage{booktabs}
\usepackage{geometry}
\geometry{a4paper, margin=2cm}

\title{熵权法（Entropy Weight）}
\author{}
\date{}

\begin{document}
\maketitle

\section{一、熵权法是干嘛的？}
给\textbf{多个指标自动算权重}。
\begin{itemize}
    \item 指标数据\textbf{差异越大 → 信息量越大 → 权重越高}
    \item 数据几乎不变 → 没信息量 → 权重越低
\end{itemize}

完全\textbf{客观}，不用人主观打分。

\section{二、举个最简单例子}
我们评价3个学生，3门课：
数学、英语、体育（都是越高越好）

原始数据矩阵：
\begin{table}[ht]
    \centering
    \begin{tabular}{|c|c|c|c|}
        \hline
        学生 & 数学 & 英语 & 体育 \\
        \hline
        A & 80 & 90 & 70 \\
        B & 90 & 85 & 80 \\
        C & 70 & 80 & 90 \\
        \hline
    \end{tabular}
    \caption{原始数据矩阵}
\end{table}

共 $n=3$ 个对象，$m=3$ 个指标

\section{三、一步一步计算（非常好懂版）}

\subsection{步骤1：数据标准化（把所有数变到0~1）}
公式：
\[ y_{ij}=\frac{x_{ij}}{\sqrt{\sum_{i=1}^n x_{ij}^2}} \]

计算：
数学平方和 = 80²+90²+70² = 19400
英语平方和 = 90²+85²+80² = 21725
体育平方和 = 70²+80²+90² = 19400

标准化后：
\begin{table}[ht]
    \centering
    \begin{tabular}{|c|c|c|c|}
        \hline
        学生 & 数学 & 英语 & 体育 \\
        \hline
        A & 80/139.28≈0.574 & 90/147.39≈0.611 & 70/139.28≈0.503 \\
        B & 90/139.28≈0.646 & 85/147.39≈0.577 & 80/139.28≈0.574 \\
        C & 70/139.28≈0.503 & 80/147.39≈0.543 & 90/139.28≈0.646 \\
        \hline
    \end{tabular}
    \caption{标准化后数据}
\end{table}

\subsection{步骤2：计算比重 pij}
\[ p_{ij}=\frac{y_{ij}}{\sum y_{ij}} \]
每一列求和，再每个数除以列和。

得到比重矩阵。

\subsection{步骤3：计算信息熵 ej}
\[ e_j = -k \sum p_{ij}\ln(p_{ij}) \]
$k=1/\ln(n)$

熵越小 → 差异越大 → 信息量越大。

\subsection{步骤4：计算差异系数 gj}
\[ g_j = 1 - e_j \]

\subsection{步骤5：计算权重 wj}
\[ w_j = \frac{g_j}{\sum g_j} \]

\section{四、例子最终权重结果（算完）}
数学：0.31
英语：0.23
体育：0.46

\textbf{解释：}
体育成绩差异最大 → 权重最高
英语成绩差异最小 → 权重最低

\section{五、熵权法总结}
\begin{enumerate}
    \item 自动算权重
    \item 数据差异越大 → 权重越高
    \item 常和 \textbf{TOPSIS} 一起用：\textbf{熵权法求权重 → TOPSIS算最终得分}
\end{enumerate}

\section{正向标准化解释}
\textbf{正向标准化 = 把所有不同大小的指标，统一缩到 [0, 1] 之间}
目的：消除单位、量纲影响，让数据可以公平计算权重。

\subsection{一、公式（正向指标：越大越好）}
对\textbf{第 j 列指标}：
\[ y_{ij} = \frac{x_{ij} - \min(x_j)}{\max(x_j) - \min(x_j)} \]
- $x_{ij}$：原始数据
- $\min(x_j)$：这一列\textbf{最小值}
- $\max(x_j)$：这一列\textbf{最大值}
- 结果：所有数 → \textbf{0 ~ 1}

\subsection{二、拿你的数据举例}
原始数学列：80，90，70
- min = 70，max = 90
- 80 → (80-70)/(90-70) = 10/20 = \textbf{0.5}
- 90 → (90-70)/20 = \textbf{1}
- 70 → 0/20 = \textbf{0}

结果：\textbf{0.5，1，0}，全部在0~1之间。

\subsection{三、为什么要标准化？}
- 数学满分100，体重50kg，身高170cm
- 单位不一样，\textbf{不能直接加权}
- 统一到0~1后，才能公平算权重

\subsection{四、代码里对应哪行？}
就是这行：
\begin{verbatim}
y = (data - x_min) / (x_max - x_min)
\end{verbatim}

\subsection{五、如果指标是越小越好？}
用\textbf{逆向标准化}：
\[ y_{ij} = \frac{\max(x_j) - x_{ij}}{\max(x_j) - \min(x_j)} \]

\end{document}